\documentclass[11pt]{article}
\usepackage[margin=1in]{geometry}
\usepackage{amsmath,amssymb,amsthm}
\usepackage[colorlinks=true,linkcolor=blue,citecolor=blue,urlcolor=blue]{hyperref}
\newtheorem{theorem}{Theorem}
\newtheorem{lemma}{Lemma}
\newtheorem{proposition}{Proposition}
\newtheorem{remark}{Remark}
\newcommand{\Fp}{\mathbb F_p}
\newcommand{\one}{\mathbf 1}
\newcommand{\E}{\mathbb E}
\newcommand{\Prob}{\mathbb P}
\newcommand{\rank}{\operatorname{rank}}
\newcommand{\nullity}{\operatorname{nullity}}
\title{Simplicity of the planted RSBM eigenvalue\\for even block sizes and internal degree at least three}
\author{Research session note}
\date{November 15, 2026}
\begin{document}
\maketitle
\begin{abstract}
For fixed integers $d_1>d_2>0$ with $d_1\ge3$, and even block size $N$, we derive high-probability simplicity of the exact planted adjacency eigenvalue $d_1-d_2$. The cross graph can be any prescribed simple $d_2$-regular bipartite graph, uniformly over that choice. The proof uses M\'esz\'aros's finite-field mixing theorem and established regular-graph contiguity, and requires no outlier gap. Exact rational elimination then gives polynomial-time recovery. The argument supplies no effective rate and does not cover $(d_1,d_2)=(2,1)$ or odd $N$ for even $d_1$. No novelty claim is made for the imported results or the planted-eigenvector identity.
\end{abstract}

\section{Statement and scope}
Fix $d_1>d_2>0$ and put $\delta=d_1-d_2$. Let $\mathcal B_{N,d_2}$ be the set of $N\times N$ zero-one matrices whose row and column sums are $d_2$. For $B\in\mathcal B_{N,d_2}$ define
\[
 A_B=\begin{pmatrix}C&B\\B^T&D\end{pmatrix},\qquad
 \sigma=\binom{\one_N}{-\one_N},
\]
where $C,D$ are independent uniform simple $d_1$-regular adjacency matrices on their respective labelled blocks. The known deterministic identity is $A_B\sigma=\delta\sigma$.

\begin{theorem}\label{thm:main}
If $d_1\ge3$, then
\[
 \lim_{\substack{N\to\infty\\N\ {\rm even}}}
 \sup_{B\in\mathcal B_{N,d_2}}
 \Prob\bigl(\nullity_{\mathbb R}(A_B-\delta I)\ge2\bigr)=0.
\]
Consequently, the same conclusion holds for an independently sampled cross graph, including the uniform simple $d_2$-regular bipartite graph in the requested RSBM.
\end{theorem}

For odd $d_1$, admissibility forces $N$ even, so this covers every admissible size. For even $d_1\ge4$, the odd-size subsequence is not covered. The only positive-degree pair $d_1>d_2$ with $d_1<3$ is $(2,1)$, also not covered. The theorem is a consequence derived here from the precise primary inputs below; it is not a verbatim theorem from those sources.

\section{The established mixing input}
Initially let $C$ be the sum of $d_1\ge3$ independent uniform perfect-matching adjacency matrices on an even set of $N$ vertices. Fix an odd prime $p$ not dividing $d_2$. For $x\in\Fp^N$, write $s(x)=\one^Tx$, and define the reference mass $q_x(r)$ as follows:
\[
 q_x(r)=\begin{cases}
  \boldsymbol 1_{\{r=d_1a\one\}},&x=a\one,\\
  p^{-N+1}\boldsymbol 1_{\{s(r)=d_1s(x)\}},&x\text{ nonconstant}.
 \end{cases}
\]
M\'esz\'aros \cite[Theorem 1.6]{meszaros} gives
\begin{equation}\label{eq:mix}
 e_N:=\sum_{x\in\Fp^N}e_x\longrightarrow0,
 \qquad e_x:=\sup_{r\in\Fp^N}|\Prob(Cx=r)-q_x(r)|.
\end{equation}
Here $p$ is fixed before $N$ tends to infinity. This is an unscaled sum of maximum-atom errors, not just convergence for a typical vector.

For completeness, the specialization of that theorem is exact. The minimal affine coset of a nonconstant vector in the additive group $\Fp$ is all of $\Fp$; for a constant vector $a\one$ it is $\{a\}$. Its additional symmetric tensor restriction lies in the subgroup generated by $a\otimes b+b\otimes a$. For odd $p$ that subgroup is all of $\Fp\otimes_{\mathbb Z}\Fp\cong\Fp$, since $2$ is invertible. Thus it imposes no additional restriction. The source mixing theorem does not require $p\nmid d_1$.

\begin{lemma}[Uniform affine bound]\label{lem:affine}
For every $b\in\Fp^N$,
\begin{align}
 \sum_x q_x(\delta x+b)&\le2,\label{eq:qaff}\\
 \sum_x\Prob(Cx=\delta x+b)&\le2+e_N.\label{eq:paff}
\end{align}
\end{lemma}
\begin{proof}
For nonconstant $x$, the reference constraint is
$d_2s(x)=s(b)$. Since $p\nmid d_2$, at most $p^{N-1}$ vectors qualify, each with mass $p^{-N+1}$. For constant $x=a\one$, a nonzero mass requires $b=d_2a\one$, which allows at most one value of $a$. This proves \eqref{eq:qaff}; summing the errors $e_x$ proves \eqref{eq:paff}.
\end{proof}

\section{The two-block first moment}
Let $C,D$ be independent copies of the matching sum, and fix $B\in\mathcal B_{N,d_2}$. All kernels in this section are over $\Fp$. Independence gives
\begin{equation}\label{eq:kernelmoment}
 \E|\ker(A_B-\delta I)|
 =\sum_{x,y}\Prob(Cx=\delta x-By)\Prob(Dy=\delta y-B^Tx).
\end{equation}
Replace the first probability by $q_x(\delta x-By)$. The total absolute error is at most
\[
 \sum_x e_x\sum_y\Prob(Dy=\delta y-B^Tx)
 \le(2+e_N)e_N
\]
by \eqref{eq:paff}. Replacing the remaining probability by $q_y(\delta y-B^Tx)$ then costs at most
\[
 \sum_y e_y\sum_xq_x(\delta x-By)\le2e_N.
\]
The combined error is at most $4e_N+e_N^2$, uniformly in $B$.

We now bound the reference sum
\[
 T_B=\sum_{x,y}q_x(\delta x-By)q_y(\delta y-B^Tx).
\]
If both vectors are nonconstant, both scalar constraints reduce to $s(x)+s(y)=0$, since $B\one=B^T\one=d_2\one$. There are at most $p^{2N-1}$ such pairs and each contributes $p^{-2N+2}$, so this part is at most $p$.

If $x=a\one$ and $y=b\one$, the constraints are exactly $a+b=0$. This contributes $p$.

Finally suppose $x=a\one$ and $y$ is nonconstant. The first constraint requires
\[
 By=-d_2a\one.
\]
For each $a$ there are at most $p^{N-r_B}$ choices, where $r_B=\rank_{\Fp}B$. Each remaining factor is at most $p^{-N+1}$. This part is at most $p^{2-r_B}$, and the opposite orientation has the same bound. We have proved
\begin{equation}\label{eq:momentbound}
 \E|\ker(A_B-\delta I)|
 \le2p+2p^{2-r_B}+4e_N+e_N^2.
\end{equation}

\begin{lemma}[Deterministic cross-matrix rank]
For every $B\in\mathcal B_{N,d_2}$ and every prime $p$,
\[
 \rank_{\Fp}B\ge \frac{N}{d_2(d_2-1)+1}.
\]
\end{lemma}
\begin{proof}
Select rows greedily with mutually disjoint supports. A selected row has $d_2$ nonzero columns; each such column belongs to $d_2-1$ other rows. Thus selecting a row excludes at most $1+d_2(d_2-1)$ rows. The selected rows have disjoint nonempty supports and are linearly independent over every field.
\end{proof}

It follows from \eqref{eq:momentbound} that, for every fixed admissible prime,
\[
 \limsup_{\substack{N\to\infty\\N\ {\rm even}}}
 \sup_B\E|\ker_{\Fp}(A_B-\delta I)|\le2p.
\]
If the real nullity is at least two, every integer minor of order $2N-1$ vanishes, so the mod-$p$ nullity is also at least two. Markov's inequality therefore gives, in the matching model,
\[
 \limsup_{\substack{N\to\infty\\N\ {\rm even}}}
 \sup_B\Prob\bigl(\nullity_{\mathbb R}(A_B-\delta I)\ge2\bigr)
 \le\frac2p.
\]
There are arbitrarily large odd primes not dividing $d_2$. Letting these fixed primes tend to infinity proves uniform high-probability simplicity in the matching model. No assertion about a prime growing with $N$ has been used.

\section{Transfer to uniform simple internal graphs}
The remaining step is established contiguity, not an equality of graph laws. The union of $d_1\ge3$ independent uniform perfect matchings on an even vertex set is mutually contiguous with the regular configuration model conditioned to have no loops. The needed result is due to Janson and to Molloy--Robalewska--Robinson--Wormald; the all-$d_1\ge3$ statement is available in \cite[Theorem 3]{mrrw}. The precise chain of models and references is also given in the proof of \cite[Theorem 1.6, Section 4, PDF p.~8]{nw}.

Finite independent products preserve contiguity. One direct formulation is that if $Q_N$ is contiguous with respect to $P_N$, then $P_N(E_N)\to0$ implies $Q_N(E_N)\to0$, uniformly for sets of sufficiently small $P_N$-probability. This follows by contradiction, choosing a violating subsequence and sets. For a product event, apply this formulation to its sections and use Markov's inequality to discard sections of non-small probability. Repeating handles the two independent internal graphs. It also permits adjoining any common independent distribution for $B$.

In particular, for any deterministic sequence $B_N\in\mathcal B_{N,d_2}$, the just-proved matching-model failure probability tends to zero and hence so does its probability under two independent loopless configuration graphs. Conditioning both on simplicity costs only a bounded factor for fixed $d_1$ and gives exactly two independent uniform simple regular graphs. Thus the conclusion holds for every deterministic sequence $B_N$ in the requested model. If the supremum in Theorem~\ref{thm:main} did not tend to zero, one could choose a violating sequence $B_N$, a contradiction. This completes the proof of Theorem~\ref{thm:main}.

\section{Exact recovery and limits of the result}
Given $A$ and $\delta$, compute the rational kernel of the integer matrix $A-\delta I$ by exact elimination. When its dimension is one, every nonzero kernel vector is a scalar multiple of $\sigma$. Normalize by its first coordinate to obtain the planted signs up to global exchange. This takes polynomial time in the bit model: fixed-degree rows have bounded Euclidean norm, so Hadamard's bound gives minors with $O(N)$ bits; standard fraction-free elimination has polynomial bit complexity. No analytic eigenvalue gap is needed.

The argument above supplies an $o(1)$ failure probability. The fixed-prime mixing limit and qualitative contiguity have not been made effective here. In particular, the theorem does not imply a polynomial error rate.

The familiar planted-eigenvector identity is already central in Barucca and the RSBM literature. Zhu--Zhu \cite[Section 1.3]{zz} explicitly calls full-parameter efficient exact recovery open. Its unqualified simplicity wording in Lemma 9.1 cannot be justified for all parameters by the cited Brito et al.\ Section 4: that argument uses even internal degree and a new-spectrum norm bound, and derives simplicity using separation from the bulk. The proof here instead uses the mixing input directly.

For $(2,1)$, a separate exact counting argument in the companion low-degree
identifiability note gives
\[
 \Prob(\text{a second admissible partition})=\Theta(N^{-2}),
 \qquad
 \Prob(\nullity_{\mathbb R}(A-I)\ge2)=\Omega(N^{-2}).
\]
Neither bound disproves high-probability real simplicity. Moreover, a certified $N=8$ example in that note has a unique sign partition and real nullity two. Thus partition uniqueness alone cannot replace the real-kernel argument. The $(2,1)$ real-kernel question and the odd-$N$, even-$d_1$ subsequence remain unresolved by this note.

\begin{thebibliography}{9}
\bibitem{meszaros}
A.~M\'esz\'aros, \emph{The distribution of sandpile groups of random regular graphs}, arXiv:1806.03736v4, 2020. Theorem~1.6 and its definitions, PDF pp.~1--4.
\url{https://arxiv.org/pdf/1806.03736v4}.

\bibitem{mrrw}
M.~S.~O.~Molloy, H.~Robalewska, R.~W.~Robinson, N.~C.~Wormald,
\emph{1-Factorizations of random regular graphs}, Random Structures \& Algorithms \textbf{10} (1997), 305--321, Theorem~3.
\href{https://onlinelibrary.wiley.com/doi/10.1002/\%28SICI\%291098-2418\%28199705\%2910\%3A3\%3C305\%3A\%3AAID-RSA1\%3E3.0.CO\%3B2-\%23}{Primary publisher page}.

\bibitem{nw}
H.~H.~Nguyen, M.~M.~Wood, \emph{Cokernels of adjacency matrices of random $r$-regular graphs}, arXiv:1806.10068v2, 2018, Section~4.
\url{https://arxiv.org/pdf/1806.10068v2}.

\bibitem{zz}
X.~Zhu, Y.~Zhu, \emph{Non-backtracking eigenvalues and eigenvectors of random regular graphs and hypergraphs}, arXiv:2312.03300v3, September 25, 2026, Sections~1.3, 3.3, 9.
\url{https://arxiv.org/html/2312.03300v3}.
\end{thebibliography}
\end{document}
